% Official Solution (Problem Sheet 6, Multiple Choice)
\begin{exercisesolution}[Solution to \cref{exc:span_subspaces_mc}]
\label{sol:span_subspaces_mc}
All four assertions are \textbf{true}. We verify each in turn:

\begin{enumerate}[label=\textbf{(\arabic*)}]
    \item \textbf{True.} The set $W = \{\lambda v \mid \lambda \in K\}$ is precisely the linear span $\Sp(v)$ of the single vector $v$. By \cref{lem:span_is_smallest}, $\Sp(S)$ is a linear subspace for any subset $S \subseteq V$. Directly: $0 = 0 \cdot v \in W$; for $\lambda_1 v, \lambda_2 v \in W$, $\lambda_1 v + \lambda_2 v = (\lambda_1 + \lambda_2) v \in W$; and for $\mu \in K$, $\mu (\lambda v) = (\mu \lambda) v \in W$.

    \item \textbf{True.} If $W \subseteq V$ is a linear subspace, then by (LSS1)--(LSS3) it contains all linear combinations of its elements, so $\Sp(W) \subseteq W$. Since $W \subseteq \Sp(W)$ always holds, we have $\Sp(W) = W$. Conversely, $\Sp(W)$ is always a linear subspace of $V$, so if $W = \Sp(W)$, then $W$ must be a linear subspace.

    \item \textbf{True.} Any vector in $\Sp(S_1 \cup S_2)$ is a finite linear combination $\sum_{i=1}^k \alpha_i u_i + \sum_{j=1}^m \beta_j v_j$ with $u_i \in S_1$ and $v_j \in S_2$. The first sum lies in $\Sp(S_1)$ and the second in $\Sp(S_2)$, so their sum lies in $\Sp(S_1) + \Sp(S_2)$, showing $\Sp(S_1 \cup S_2) \subseteq \Sp(S_1) + \Sp(S_2)$. Conversely, $S_1 \subseteq S_1 \cup S_2 \implies \Sp(S_1) \subseteq \Sp(S_1 \cup S_2)$ and $S_2 \subseteq S_1 \cup S_2 \implies \Sp(S_2) \subseteq \Sp(S_1 \cup S_2)$. Since $\Sp(S_1 \cup S_2)$ is a subspace containing both $\Sp(S_1)$ and $\Sp(S_2)$, it contains their sum $\Sp(S_1) + \Sp(S_2)$.

    \item \textbf{True.} Since $S_1 \cap S_2 \subseteq S_1$, we have $\Sp(S_1 \cap S_2) \subseteq \Sp(S_1)$. Similarly, $S_1 \cap S_2 \subseteq S_2$ implies $\Sp(S_1 \cap S_2) \subseteq \Sp(S_2)$. Therefore, $\Sp(S_1 \cap S_2) \subseteq \Sp(S_1) \cap \Sp(S_2)$.
\end{enumerate}
\exinfo{Official Solution (Problem Sheet 6, Multiple Choice). Verifies subspace span properties and intersection inclusions.}
\end{exercisesolution}

% Official Solution (Problem Sheet 6, Exercise 2)
\begin{exercisesolution}[Solution to \cref{exc:subspaces_two_dimensions}]
\label{sol:subspaces_two_dimensions}
\begin{enumerate}[label=\textbf{(\alph*)}]
    \item \textbf{Proof that $\Sp(v, w) = \mathbb{R}^2$:}
    Let $v = (v_1, v_2) \ne (0, 0)$ and $w = (w_1, w_2) \ne (0, 0)$ with $w \ne \alpha v$ for all $\alpha \in \mathbb{R}$.
    Since $w$ is not a scalar multiple of $v$, the quantity $\Delta := v_1 w_2 - v_2 w_1 \ne 0$. (If $\Delta = 0$, then with $v_1 \ne 0$ we would have $w_2 = \frac{v_2 w_1}{v_1}$, so $w = \frac{w_1}{v_1} v$, a contradiction.)
    
    Given any arbitrary vector $u = (u_1, u_2) \in \mathbb{R}^2$, we seek scalars $c_1, c_2 \in \mathbb{R}$ such that $c_1 v + c_2 w = u$, which in coordinates reads:
    \[
    \begin{cases}
    c_1 v_1 + c_2 w_1 = u_1 \\
    c_1 v_2 + c_2 w_2 = u_2
    \end{cases}
    \]
    Since $\Delta \ne 0$, this $2 \times 2$ system has the unique solution:
    \[
    c_1 = \frac{u_1 w_2 - u_2 w_1}{\Delta}, \qquad c_2 = \frac{v_1 u_2 - v_2 u_1}{\Delta}.
    \]
    Thus $u = c_1 v + c_2 w \in \Sp(v, w)$, proving $\Sp(v, w) = \mathbb{R}^2$.

    \item \textbf{Classification of all Subspaces of $\mathbb{R}^2$:}
    Let $U \subseteq \mathbb{R}^2$ be an arbitrary linear subspace.
    \begin{itemize}
        \item If $U$ contains only the zero vector, then $U = \{(0, 0)\}$.
        \item If $U \ne \{(0, 0)\}$, there exists some non-zero vector $v \in U$. By (LSS3), $\Sp(v) \subseteq U$.
        \begin{itemize}
            \item If $U = \Sp(v)$, then $U$ is a line through the origin.
            \item If $U \ne \Sp(v)$, there exists a vector $w \in U$ such that $w \notin \Sp(v)$. Then $w \ne \alpha v$ for all $\alpha \in \mathbb{R}$. By part \textbf{(a)}, $\Sp(v, w) = \mathbb{R}^2$. Since $v, w \in U$ and $U$ is a subspace, we have $\mathbb{R}^2 = \Sp(v, w) \subseteq U \subseteq \mathbb{R}^2$, which forces $U = \mathbb{R}^2$.
        \end{itemize}
    \end{itemize}
    Thus, the only subspaces of $\mathbb{R}^2$ are $\{(0, 0)\}$, lines through the origin $\Sp(u)$ ($u \ne 0$), and $\mathbb{R}^2$ itself.
\end{enumerate}
\exinfo{Official Solution (Problem Sheet 6, Exercise 2). Proves two non-collinear vectors span $\mathbb{R}^2$ and classifies all subspaces of $\mathbb{R}^2$.}
\end{exercisesolution}

% Solution: Gemini / Official Hybrid (Problem Sheet 6, Exercise 1)
\begin{exercisesolution}[Solution to \cref{exc:polynomial_linear_combination_span}]
\label{sol:polynomial_linear_combination_span}
\begin{enumerate}[label=\textbf{(\alph*)}]
    \item \textbf{Expressing $q(x)$ as a Linear Combination:}
    We seek scalars $\alpha_1, \alpha_2, \alpha_3, \alpha_4 \in \mathbb{R}$ such that
    \[
    \alpha_1 (x^3 + x^2) + \alpha_2 (x^2 - 2x - 4) + \alpha_3 (3x + 4) + \alpha_4 (2x + 3) = 2x^3 + 3x^2 - 1.
    \]
    Equating coefficients of $x^3, x^2, x^1, x^0$:
    \begin{align*}
    [x^3]: & \quad \alpha_1 = 2, \\
    [x^2]: & \quad \alpha_1 + \alpha_2 = 3 \implies 2 + \alpha_2 = 3 \implies \alpha_2 = 1, \\
    [x^1]: & \quad -2\alpha_2 + 3\alpha_3 + 2\alpha_4 = 0 \implies -2(1) + 3\alpha_3 + 2\alpha_4 = 0 \implies 3\alpha_3 + 2\alpha_4 = 2, \\
    [x^0]: & \quad -4\alpha_2 + 4\alpha_3 + 3\alpha_4 = -1 \implies -4(1) + 4\alpha_3 + 3\alpha_4 = -1 \implies 4\alpha_3 + 3\alpha_4 = 3.
    \end{align*}
    Solving the $2 \times 2$ system for $\alpha_3, \alpha_4$:
    \[
    \begin{cases}
    3\alpha_3 + 2\alpha_4 = 2 \\
    4\alpha_3 + 3\alpha_4 = 3
    \end{cases}
    \implies
    \alpha_3 = 3(2) - 2(3) = 0, \quad \alpha_4 = 3(3) - 4(2) = 1.
    \]
    Thus, the linear combination is:
    \[
    q(x) = 2 p_1(x) + 1 p_2(x) + 0 p_3(x) + 1 p_4(x) = 2 p_1(x) + p_2(x) + p_4(x).
    \]

    \item \textbf{Determination of $\Sp(p_1, p_2, p_3, p_4)$:}
    Since each $p_i(x)$ has degree at most $3$, we have $\Sp(p_1, p_2, p_3, p_4) \subseteq \mathbb{R}[x]_3$. We show that the span contains the standard monomial basis $\{1, x, x^2, x^3\}$:
    \begin{itemize}
        \item $3 p_3(x) - 4 p_4(x) = 3(3x + 4) - 4(2x + 3) = x \in \Sp(p_1, p_2, p_3, p_4)$.
        \item $p_3(x) - 3x = 4 \implies 1 = \frac{1}{4}(p_3(x) - 3x) \in \Sp(p_1, p_2, p_3, p_4)$.
        \item $x^2 = p_2(x) + 2x + 4 \in \Sp(p_1, p_2, p_3, p_4)$.
        \item $x^3 = p_1(x) - x^2 \in \Sp(p_1, p_2, p_3, p_4)$.
    \end{itemize}
    Since $\Sp(1, x, x^2, x^3) = \mathbb{R}[x]_3$, we conclude that
    \[
    \Sp(p_1, p_2, p_3, p_4) = \mathbb{R}[x]_3.
    \]
\end{enumerate}
\exinfo{Hybrid Solution (Gemini 3.7 Flash / Official Problem Sheet 6, Exercise 1). Solves the linear system for polynomial coordinates and proves the polynomials span $\mathbb{R}[x]_3$.}
\end{exercisesolution}

% Official Solution (Problem Sheet 6, Exercise 3)
\begin{exercisesolution}[Solution to \cref{exc:matrix_subspace_vector_constraint}]
\label{sol:matrix_subspace_vector_constraint}
The subset $U = \{A \in \M_{m \times n}(K) \mid A \cdot x = b\}$ is a linear subspace of $\M_{m \times n}(K)$ if and only if $b = 0 \in K^m$ (for any $x \in K^n$).

\begin{itemize}
    \item \textbf{Necessity of $b = 0$:}
    If $U$ is a linear subspace, it must contain the zero matrix $0 \in \M_{m \times n}(K)$ by (LSS1). Since $0 \cdot x = 0 \in K^m$ for any vector $x$, the condition $0 \in U$ forces $b = 0 \cdot x = 0$.

    \item \textbf{Sufficiency when $b = 0$:}
    Assume $b = 0$. Then $U = \{A \in \M_{m \times n}(K) \mid A \cdot x = 0\}$. We check the subspace criteria:
    \begin{enumerate}[label=\textbf{(\alph*)}]
        \item The zero matrix satisfies $0 \cdot x = 0$, so $0 \in U$ and $U \ne \emptyset$.
        \item Let $A, B \in U$ and $\alpha, \beta \in K$. Then by the linearity of matrix-vector multiplication (\cref{lem:matrix_vector_linearity}):
        \[
        (\alpha A + \beta B) \cdot x = \alpha (A \cdot x) + \beta (B \cdot x) = \alpha \cdot 0 + \beta \cdot 0 = 0.
        \]
        Thus $\alpha A + \beta B \in U$.
    \end{enumerate}
\end{itemize}
Therefore, $U$ is a linear subspace if and only if $b = 0$.
\exinfo{Official Solution (Problem Sheet 6, Exercise 3). Proves $Ax = b$ subspace criterion requires $b=0$.}
\end{exercisesolution}

% Official Solution (Problem Sheet 6, Exercise 4)
\begin{exercisesolution}[Solution to \cref{exc:polynomial_space_infinite_dim}]
\label{sol:polynomial_space_infinite_dim}
Suppose for contradiction that the vector space $K[x]$ is finite-dimensional. Then by \cref{def:finite_dimensional}, there exists a finite generating set $S = \{f_1, \dots, f_k\} \subset K[x]$ such that $\Sp(S) = K[x]$.

Each polynomial $f_i(x)$ has a well-defined non-negative degree $\deg(f_i)$ (or degree $-\infty$ if $f_i = 0$). Since the set $S$ is finite, the maximum
\[
d := \max_{1 \le i \le k} \deg(f_i) \in \mathbb{Z}_{\ge 0} \cup \{-\infty\}
\]
exists. For any polynomial $g(x) \in \Sp(f_1, \dots, f_k)$, $g(x)$ is a finite linear combination $\sum_{i=1}^k c_i f_i(x)$ with $c_i \in K$. By the properties of polynomial degrees,
\[
\deg(g) \le \max_{1 \le i \le k} \deg(f_i) = d.
\]
Consequently, every element of $\Sp(S)$ has degree at most $d$, so $\Sp(S) \subseteq K[x]_d$.
However, the monomial $x^{d+1}$ belongs to $K[x]$ and has $\deg(x^{d+1}) = d + 1 > d$, so $x^{d+1} \notin \Sp(S)$.
This contradicts the assumption that $\Sp(S) = K[x]$. Therefore, no finite set can span $K[x]$, and $K[x]$ is not finite-dimensional.
\exinfo{Official Solution (Problem Sheet 6, Exercise 4). Demonstrates $K[x]$ is infinite-dimensional via the degree bound.}
\end{exercisesolution}

% Official Solution (Problem Sheet 6, Exercise 5)
\begin{exercisesolution}[Solution to \cref{exc:generating_sets_sequences}]
\label{sol:generating_sets_sequences}
For each $i \in \mathbb{Z}_{\ge 0}$, let $e_i := (0, \dots, 0, 1, 0, \dots) \in K^\infty$ denote the sequence with $1$ in position $i$ and $0$ in all other positions.

\begin{enumerate}[label=\textbf{(\alph*)}]
    \item \textbf{Smallest Generating Set for $K_0^\infty$:}
    Consider the set $E := \{e_i \mid i \ge 0\} \subseteq K_0^\infty$.
    \begin{itemize}
        \item \textbf{$E$ spans $K_0^\infty$:} Any sequence $a = (a_0, a_1, \dots) \in K_0^\infty$ has only finitely many non-zero terms, say $a_n = 0$ for all $n > N$. Then $a = \sum_{i=0}^N a_i e_i$, which is a finite linear combination of elements of $E$. Hence $\Sp(E) = K_0^\infty$.
        \item \textbf{Minimality:} If $E' \subsetneq E$, there exists some index $k \ge 0$ such that $e_k \notin E'$. Any finite linear combination of elements in $E'$ has $k$-th coordinate equal to $0$, because every $e_j \in E'$ ($j \ne k$) has $(e_j)_k = 0$. Hence $e_k \notin \Sp(E')$, so $\Sp(E') \ne K_0^\infty$.
    \end{itemize}
    Thus, $E = \{e_i \mid i \ge 0\}$ is the unique minimal (under inclusion, up to scalar multiples) generating set of $K_0^\infty$.

    \item \textbf{Smallest Generating Set for $K_{\mathrm{cst}}^\infty$:}
    Let $\mathbf{1} := (1, 1, 1, \dots) \in K_{\mathrm{cst}}^\infty$ be the constant sequence of ones. Consider the set $S := \{e_i \mid i \ge 0\} \cup \{\mathbf{1}\} \subseteq K_{\mathrm{cst}}^\infty$.
    \begin{itemize}
        \item \textbf{$S$ spans $K_{\mathrm{cst}}^\infty$:} Let $a = (a_0, a_1, \dots) \in K_{\mathrm{cst}}^\infty$, so there exist $c \in K$ and $N \ge 0$ such that $a_n = c$ for all $n > N$. The sequence $a - c \cdot \mathbf{1}$ has $n$-th entry $0$ for all $n > N$, so $a - c \cdot \mathbf{1} \in K_0^\infty$. Writing $a - c \cdot \mathbf{1} = \sum_{i=0}^N (a_i - c) e_i$, we obtain
        \[
        a = \sum_{i=0}^N (a_i - c) e_i + c \cdot \mathbf{1} \in \Sp(S).
        \]
        Thus $\Sp(S) = K_{\mathrm{cst}}^\infty$.
        \item \textbf{Minimality:}
        \begin{itemize}
            \item If $\mathbf{1}$ is omitted, $\Sp(E) = K_0^\infty \subsetneq K_{\mathrm{cst}}^\infty$, because $\mathbf{1} \notin K_0^\infty$.
            \item If some $e_k$ is omitted from $S$, consider $\Sp(S \setminus \{e_k\})$. Any element in this span has the form $\sum_{j \ne k} \alpha_j e_j + \beta \mathbf{1}$. If this sequence equals $e_k$, its tail entries must be $0$, which forces $\beta = 0$. But with $\beta = 0$, its $k$-th coordinate is $0 \ne 1$. Hence $e_k \notin \Sp(S \setminus \{e_k\})$.
        \end{itemize}
    \end{itemize}
    Thus, $S = \{e_i \mid i \ge 0\} \cup \{\mathbf{1}\}$ is the minimal generating set of $K_{\mathrm{cst}}^\infty$.
\end{enumerate}
\exinfo{Official Solution (Problem Sheet 6, Exercise 5). Constructs minimal generating sets for finite-support and eventually-constant sequence spaces.}
\end{exercisesolution}
